\section{Methods}

\subsection{Pipeline Architecture}

GlycoSMILES2BAP operates through four sequential modules (Figure~1):
\begin{enumerate}
\item \textbf{Input Parser:} Converts glycan notation to Abstract Syntax Tree (AST)
\item \textbf{CCD Mapper:} Assigns Chemical Component Dictionary codes
\item \textbf{BAP Generator:} Creates bondedAtomPairs specifications
\item \textbf{Output Formatter:} Produces AF3-compatible JSON input
\end{enumerate}

\subsection{Input Parsing Module}

The pipeline accepts three input formats commonly used in glycobiology:
\begin{itemize}
\item \textbf{IUPAC-condensed:} The standard notation used in publications (e.g., \texttt{Gal(b1-4)GlcNAc})
\item \textbf{WURCS:} Web3 Unique Representation of Carbohydrate Structures
\item \textbf{GlycoCT:} A comprehensive XML-based glycan format
\end{itemize}

The parser module leverages existing glycan parsing tools (GlyLES, GlycanFormatConverter) to convert input strings into a standardized AST representation capturing: monosaccharide identity, anomeric configuration ($\alpha$/$\beta$), absolute configuration (D/L), ring type, modifications, and glycosidic linkage positions.

\subsection{CCD Mapper Module}

The CCD (Chemical Component Dictionary) provides standardized residue identifiers used by the Protein Data Bank. Our mapper supports 28+ monosaccharide configurations.

\begin{table}[h]
\centering
\caption{Key CCD Mappings}
\label{tab:ccd}
\begin{tabular}{lllll}
\toprule
Monosaccharide & Anomer & Config & CCD Code & Anomeric C \\
\midrule
GlcNAc & $\beta$ & D & NAG & C1 \\
GlcNAc & $\alpha$ & D & A2G & C1 \\
Man & $\alpha$ & D & MAN & C1 \\
Man & $\beta$ & D & BMA & C1 \\
Gal & $\beta$ & D & GAL & C1 \\
Gal & $\alpha$ & D & GLA & C1 \\
Fuc & $\alpha$ & L & FUC & C1 \\
Neu5Ac & $\alpha$ & D & SIA & C2 \\
Neu5Gc & $\alpha$ & D & NGC & C2 \\
Xyl & $\beta$ & D & XYS & C1 \\
\bottomrule
\end{tabular}
\end{table}

Key design decisions:
\begin{enumerate}
\item \textbf{Case-insensitive matching:} All monosaccharide names are normalized to lowercase for robust lookup
\item \textbf{Anomeric position tracking:} Sialic acids (Neu5Ac, Neu5Gc, KDN) use C2 as anomeric carbon (ketoses), while aldoses use C1
\item \textbf{Ring oxygen positions:} Pentoses (Xyl, Ara, Rib) use O4; hexoses use O5; sialic acids use O6
\end{enumerate}

\subsection{BAP Generator Module}

The bondedAtomPairs (BAP) generator converts glycan topology into AF3-compatible bond specifications. Each glycosidic linkage is represented as an explicit atom pair.

\textbf{Branch Topology Resolution.} For glycans with complex branched topologies such as bi-antennary N-glycans (e.g., G2 structures with Man$\alpha$1-3[Man$\alpha$1-6]Man branches), the BAP generator employs a \textbf{depth-first search (DFS)} algorithm to traverse the Abstract Syntax Tree (AST) and systematically generate all bond pairs. The DFS traversal proceeds from the reducing end toward the non-reducing ends, ensuring that each branch is fully processed before backtracking to explore sibling branches. This approach guarantees complete coverage of all glycosidic linkages while maintaining the correct donor-acceptor residue indexing required for AF3 input. The traversal algorithm handles multiple branching points by recursively descending into each child branch, generating BAP entries for each glycosidic bond encountered before returning to the parent node. This ensures that complex multi-antennary structures are correctly represented with all inter-residue bonds explicitly specified.

\textbf{Algorithmic complexity.} The DFS-based BAP generation operates in $O(n)$ time complexity, where $n$ is the number of monosaccharide residues, as each residue and its associated linkage is visited exactly once during the traversal. The space complexity is $O(d)$ where $d$ is the maximum depth of the glycan tree structure, corresponding to the DFS stack depth required for traversal.

For example, the Gal($\beta$1-4)GlcNAc linkage generates:
\begin{verbatim}
{
  "residue1": 1,
  "atom1": "C1",
  "residue2": 2,
  "atom2": "O4",
  "order": 1
}
\end{verbatim}

\subsection{Error Handling Strategy}

The pipeline implements validation at each processing stage:
\begin{enumerate}
\item \textbf{Parser validation:} Syntax checking for IUPAC-condensed and WURCS input strings
\item \textbf{CCD mapper validation:} Hierarchical fallback system for unknown monosaccharides
\item \textbf{BAP generator validation:} Topology consistency checks for valid donor-acceptor indices
\end{enumerate}

\subsection{Software Dependencies}

GlycoSMILES2BAP requires: Python 3.8+, glyles library, glypy, and standard libraries (json, re, typing). The tool is compatible with Linux, macOS, and Windows.

\subsection{Computational Environment}

All benchmark performance measurements were conducted on a standard desktop computing environment: Intel Core i7-10700 CPU @ 2.90GHz (8 cores, single-threaded execution), 16GB DDR4 RAM, running Ubuntu 20.04 LTS. Processing time measurements represent wall-clock time for complete pipeline execution (parsing, CCD mapping, and BAP generation), averaged over 3 runs per structure. No GPU acceleration was utilized, as the computational workload is entirely CPU-bound and does not benefit from parallel processing for individual glycan structures.

\subsection{Evaluation Metrics}
\label{sec:metrics}

To rigorously assess pipeline performance, we defined three complementary accuracy metrics. \textbf{Importantly, all metrics are calculated at the level of individual structural features (per monosaccharide residue or per glycosidic linkage), not at the level of complete molecular structures.} This granular approach enables precise identification of specific stereochemistry preservation capabilities.

\textbf{Epimer accuracy} measures correct preservation of monosaccharide stereochemistry at all chiral centers. For each monosaccharide residue $i$ in the glycan structure, the epimer status is considered correct if the generated CCD code matches the expected CCD code for that specific (monosaccharide, anomer, absolute configuration) combination. The epimer accuracy is calculated as:

\begin{equation}
\text{Epimer Accuracy} = \frac{\text{Number of correctly identified monosaccharide residues}}{\text{Total number of monosaccharide residues in benchmark}}
\end{equation}

For example, if a benchmark glycan contains 4 monosaccharide residues and all 4 are correctly mapped to their corresponding CCD codes, the epimer accuracy for that structure is 100\%. A residue is counted as incorrect if Gal is misidentified as Glc (epimer error), or if GlcNAc is misidentified as Glc (modification error).

\textbf{Anomeric accuracy} measures correct preservation of the anomeric configuration ($\alpha$ or $\beta$) for each glycosidic linkage. For each linkage $j$, the anomeric configuration is correct if the donor anomeric carbon designation (C1 for aldoses, C2 for ketoses/sialic acids) matches the expected configuration. The anomeric accuracy is calculated as:

\begin{equation}
\text{Anomeric Accuracy} = \frac{\text{Number of linkages with correct anomer}}{\text{Total number of glycosidic linkages in benchmark}}
\end{equation}

For a linear tetrasaccharide with 3 glycosidic linkages, if 3 linkages have correct $\alpha$/$\beta$ designation, the anomeric accuracy is 100\%.

\textbf{Linkage accuracy} measures correct specification of both donor and acceptor atom positions for each glycosidic linkage. A linkage is considered correctly specified if: (1) the donor anomeric carbon position matches the expected position, (2) the acceptor oxygen position matches the expected linkage position, and (3) the residue indices correctly reflect the glycan topology. The linkage accuracy is calculated as:

\begin{equation}
\text{Linkage Accuracy} = \frac{\text{Number of linkages with correct BAP specification}}{\text{Total number of glycosidic linkages in benchmark}}
\end{equation}

\textbf{Sample size for calculations:} In our benchmark of 50 glycan structures containing a total of 178 monosaccharide residues and 128 glycosidic linkages, each metric was calculated by aggregating across all residues or linkages, respectively. This approach provides a more stable and informative measure of pipeline performance than whole-structure accuracy, which would be overly stringent (a single error would render an entire structure ``incorrect'')
